% ===========================================================================
% 06-Conclusions.tex.
% Terminology, load-bearing: held-out treated WELL, trained CONDITION.
% "treated well" is the glossary's form (\glentry, 00-HowToRead:267) and wins
% over "treatment well" everywhere. Never "pairing"; the thesis declines that
% reading.
% ===========================================================================

\chapter{Conclusions}
\label{sec:conclusions}
\framedecl{ER}

\begin{question}
Does a cell-sentence language model fine-tuned on this atlas learn what a
particular drug does to a particular cell? The instrument was a chain built in
order: a metric calibrated before it was trusted, controls a drug-blind
predictor cannot pass, and a bar that needs no model at all. The bar was not
cleared.
\end{question}

\noindent
Prompt sensitivity to trained drug names is established; biological transfer is
\notestablished; and nothing built here beats a training-only per-drug lookup on
common support. The win condition was fixed in advance
(\cref{sec:introduction}): a conditional model must beat that lookup, because
only then has its access to the untreated cell supplied something a dictionary
of per-drug averages lacks. A split-half reference cannot stand in for it, since
it cannot separate a drug's characteristic response, which transfers, from this
context's departure from it. That verdict ends a chain, and no isolated point
estimate may leapfrog it.

% ---------------------------------------------------------------------------
\section{Summary of the argument}

\paragraph{The metric had to be settled first} Four of five candidate metrics failed the calibration audit, and the field's standard one is saturated outright by a predictor carrying no drug information: on $\DEdr$ such a predictor scores $0.9999$. Only $\NIR$ registered known differences in quality,
at a dynamic range fraction of \ci{+0.635}{+0.604}{+0.663}, so it is the only
metric this thesis grades anything on. That is a result before it is an
apparatus choice, and it decides what every later number is allowed to mean.

\paragraph{On that metric the expression-frame checkpoint is drug-blind} It
scores $0.498$ against a within-well split-half reference of $0.576$. On the
stratum where the ranking is provably winnable, a control copying the untreated
cell matches it outright, $0.768$ against $0.766$, a paired contrast of
\ci{+0.002}{-0.026}{+0.028}. Three instruments that fail in different ways agree:
the scramble, the control copy, and a within-condition regression of score on how
unlike the substituted drug is, which returns \ci{+0.0055}{-0.0472}{+0.0582} once
each condition is compared only with itself. That last one matters because its
pooled version returns $+0.0986$ and looks like an effect; the difference between
the two is condition-level structure, not the drug.

\paragraph{The drug is present, and barely consulted} The null is not a failure
to measure. The drug label is decodable from the residual stream at $82\%$
against an $8\%$ chance rate, so it survives into the computation. Ablating the
purified drug subspace costs the output $3.6$ to $18.6$ times a matched-dimension random null, and $9$ to $37\%$ of what ablating cell line costs at six of seven depths ($98\%$ at layer $12$, where both costs collapse). The label is
read, and read weakly, which is a different failure from absence and needs a
different repair.

% ---------------------------------------------------------------------------
\section{Prompt sensitivity on trained drugs}

If the drug is in the data and in the activations, the next suspect is the
training target. Two drugs in one context share $83\%$ of their standard target's top-$200$ gene positions; the residual encoding drops that to $40\%$. Aggregating the
target without re-referencing it does not work: of three constructions tried,
none separates from a random-partner scramble, at $+0.014$, $-0.001$ and
$-0.010$, and the optimal-transport arm's within-condition slope is
\ci{-0.0008}{-0.0305}{+0.0290}. No positive slope is established at this
precision; the interval does not establish equivalence to zero.

Re-encoding does. On conditions the checkpoint trained on, the residual arm
separates from its own scramble against the geometry-defined \texttt{orth}
partner by $+0.0898$, excluding zero clustered on cell line and well,
\ci{+0.0898}{+0.0382}{+0.1415}, and on drug and well,
\ci{+0.0898}{+0.0291}{+0.1506}.

\paragraph{The comparator is part of the measurement} A prompt replacement is
not automatically a null. Pooled over all $\num{1394}$ scored conditions, the scrambled arm itself scores $0.550$, $0.501$ and $0.423$ as the substituted drug moves from aligned through orthogonal to
anti-aligned, so a dissimilar lie inflates the apparent benefit of the correct
name. \texttt{Orth} is the one comparator whose own score sits at chance, and it
is the only one whose geometry holds still across the splits the next section
compares.

\paragraph{The gain identifies a package, not a mechanism} Residual encoding
changes aggregation, control referencing, generic subtraction, sign coding and
sequence length at once, and the factorial comparison that would separate them
was not run. Nor is it uniquely responsible: at a matched $1400$-token budget the
optimal-transport arm also clears zero against its orthogonal comparator,
\ci{+0.0292}{+0.0144}{+0.0441}, where its $600$-token estimate spanned zero. A
generation cap can manufacture a null, so any claim about the residual advantage
carries its budget with it.

% ---------------------------------------------------------------------------
\section{Transfer}

One question is whether the prediction moves when the drug named in the prompt
changes. The other is whether it moves \emph{toward} that drug's response. The
trained-condition effect settles the first for drugs represented in fine-tuning
and leaves the second open.

Conditions in held-out treated wells score $+0.0332$, excluding zero clustered on
cell line and well, \ci{+0.0332}{+0.0069}{+0.0594}, but not when clustered on drug,
\ci{+0.0332}{-0.0031}{+0.0695} --- and drug is the axis a transferred residual
has to generalise over. That split is largely recombination in any case: $30$
of its $34$ drugs and $39$ of its $42$ cell lines also appear in the sampled
\texttt{train} evaluation, and $82\%$ of its conditions pair a drug and cell
line present in that sample. Absence from the sample is not absence from
fine-tuning: all four other drugs have training-only lookup support. The
separate \texttt{unseen\_drug} split contains the $10$ drugs withheld from
fine-tuning (\cref{sec:methods-holdout}). The model's own
score covers chance on both held-out splits.

\paragraph{What survives the holdout is sensitivity, not accuracy} The
within-condition slope is positive on all three splits and under both
clusterings, \ci{+0.361}{+0.249}{+0.472}, \ci{+0.321}{+0.207}{+0.434} and
\ci{+0.324}{+0.146}{+0.502}, on an instrument that needs no comparator and whose
null is fixed by construction. On the same instrument the single-cell and
optimal-transport arms return nothing. But every point on that line is a
\emph{wrong} prompt, so it shows the model tells wrong answers apart and never
that it produces a right one; and on the unseen-drug split the drugs it names are
overwhelmingly training drugs. The map is coarse and trained-drug-to-residual.

\begin{scopelimit}[discharges the limit carried from \cref{sec:limitations}]
The \texttt{unseen\_drug} arm is read here as an instrument control. Its
estimated model-minus-neutral-scramble gap is compatible with zero and a range
of small effects under both reported clusterings. This does not establish
equivalence to zero or the absence of unseen-drug transfer, and the prompt's
mechanism field limits the claim to held-out drug identities.
\end{scopelimit}

% ---------------------------------------------------------------------------
\section{The lookup baseline}

A per-drug lookup needs no checkpoint, no prompt and no cell
sentence.\seealso{\cref{sec:res-lookup} scores the full ladder through one
pipeline} It predicts a drug's residual as its mean residual in other cell
lines, fitted from training conditions only. On the $\num{1192}$ conditions where
both are defined it scores $0.913$ against the model's $0.549$, with the
within-well reference at $0.857$. Restricted to a single other cell line it still
scores $0.822$, so the win is not an artefact of averaging away noise.

The reading was fixed before the comparison ran: above chance and below the
lookup means the model recovers something of a drug's average residual and adds
no tailoring to the cell line in front of it. That is where it landed. This is
DrEval's finding with a transcriptional vector in place of a scalar potency, and
it sets the bar a conditional architecture has to clear to earn its complexity.

% ---------------------------------------------------------------------------
\section{Unseen drugs}

No lookup exists for a drug absent from training, so the last gate asks which
shared property could stand in for one. Three channels were tested against a
plate-matched random null, with partners restricted to training conditions.

Protein target clears the pre-set relevance margin clustered on cell line and
well, \ci{+0.1351}{+0.0745}{+0.1958}. Clustered on the drug --- the axis an
unseen-drug claim rides on --- it clears the margin only against the weaker
count-matched null and falls short against both estimand-correct ones. Mechanism
excludes zero on the cell-line axis and misses the margin by $0.0021$; chemistry
straddles it. All three exclude zero, so the two that fail are inconclusive
rather than equivalent to their nulls, and none of the three may be written as
closed.

The lead rests on $18$ target-annotated drugs, two of them held out. Protein
target carries transcriptional signal on that support, and a Leave-Drugs-Out
result is \notestablished by it. Which span of the prompt the model actually
reads is unresolved (\cref{sec:res-field}): the two spans were swapped separately, but the results cannot be told apart, because each was scored against a comparator on a different side of chance (partner $\NIR$ $0.434$ and $0.520$; \cref{tab:fieldattribution}, Panel B), so a channel-conditioned prompt would add a field
to an instruction whose existing fields are of unmeasured value. The next
experiment needs a target-annotated panel broad enough to cluster on, preferably
with held-out mechanism families, and a prompt whose spans have been told
apart.

% ---------------------------------------------------------------------------
\section{The role of the controls}

In a literature where uninformative baselines score well, a negative result is
worth what its controls are worth. Each surviving claim here rests on a control
added after a specific failure: plate-matched comparison sets after a
zero-information baseline improved when plate was free to vary; a swap-distance
sweep after a similar-drug scramble proved weak; a removal gate before any
ablation null; matched-norm injection before any swap was interpreted; condition
intercepts after a pooled slope returned an effect that was not there. Intervals
were clustered on the units that repeat --- cell line, drug and treated well ---
and never on cells.

Claims that failed those controls were withdrawn and the withdrawals recorded
rather than quietly dropped. That record is why the remaining claims can be read
at the strength stated here, and no stronger.

\begin{claim}
On the $\num{1192}$ conditions where both are defined, a training-only per-drug
lookup scores $0.913$ against the model's $0.549$, with the within-well reference
at $0.857$; from a single other cell line it still scores $0.822$. Residual
re-encoding produces a real trained-condition effect against the \texttt{orth}
partner, excluding zero on both clustering axes, and establishes no transfer:
held-out treated wells give an interval that excludes zero on the cell-line axis
and spans zero on the drug axis, and the model's own score covers chance on both
held-out splits. What survives the holdout is the comparator-free slope, which is
sensitivity to the drug named and not accuracy about the drug given. No drug
effect is \shownzero by any instrument here. Of the three channels open to a drug
with no lookup, protein target is live on the cell-line axis and inconclusive on
the drug axis, mechanism and chemistry are inconclusive, and none is closed.
\end{claim}
